\documentclass[10pt,a4paper]{article}
\usepackage[margin=1.6cm, top=2cm, bottom=2cm]{geometry}
\usepackage[T1]{fontenc}
\usepackage[utf8]{inputenc}
\usepackage{listings}
\usepackage{xcolor}
\usepackage{booktabs}
\usepackage{tabularx}
\usepackage{multicol}
\usepackage{parskip}
\usepackage{titlesec}
\usepackage{enumitem}
\usepackage{fancyhdr}
\usepackage{tcolorbox}
\usepackage{colortbl}
\usepackage{array}
\usepackage{amsmath}

\tcbuselibrary{skins,breakable}

\pagestyle{fancy}
\fancyhf{}
\rhead{\small\textcolor{sqlpurple}{\textbf{SQL Samenvatting}} --- IEE Databank --- Databeheer 2025--2026}
\lhead{\small Louic Cooreman}
\rfoot{\small\thepage}
\renewcommand{\headrulewidth}{0.4pt}

\definecolor{sqlpurple}{RGB}{136,60,180}
\definecolor{sqlblue}{RGB}{50,120,200}
\definecolor{sqlgreen}{RGB}{30,140,80}
\definecolor{sqlorange}{RGB}{200,110,10}
\definecolor{sqlgray}{RGB}{110,110,110}
\definecolor{sqlred}{RGB}{170,30,30}
\definecolor{codbg}{RGB}{246,243,255}
\definecolor{resbg}{RGB}{241,251,241}
\definecolor{tblhead}{RGB}{80,30,120}
\definecolor{notecolor}{RGB}{255,249,220}
\definecolor{warncolor}{RGB}{255,236,236}
\definecolor{rowalt}{RGB}{248,244,255}

\lstdefinestyle{sql}{
  language=SQL,
  basicstyle=\ttfamily\footnotesize,
  keywordstyle=\bfseries\color{sqlpurple},
  stringstyle=\color{sqlorange},
  commentstyle=\itshape\color{sqlgray},
  backgroundcolor=\color{codbg},
  frame=single, framerule=0.5pt,
  rulecolor=\color{sqlpurple!35},
  breaklines=true, showstringspaces=false,
  tabsize=2, xleftmargin=5pt, xrightmargin=5pt,
  morekeywords={INNER,LEFT,RIGHT,JOIN,ON,AS,REFERENCES,AUTO_INCREMENT,
    DATETIME,VARCHAR,INT,CHAR,DOUBLE,FLOAT,TEXT,YEAR,DATE,NOT,NULL,
    DEFAULT,UNIQUE,PRIMARY,KEY,FOREIGN,CONSTRAINT,DROP,IF,
    EXISTS,USE,CREATE,DATABASE,TABLE,INSERT,INTO,VALUES,UPDATE,SET,
    DELETE,FROM,WHERE,GROUP,BY,HAVING,ORDER,SELECT,DISTINCT,COUNT,
    AVG,SUM,MIN,MAX,AND,OR,IN,BETWEEN,LIKE,DESC,ASC,NOW,CONCAT,
    IS,PRECISION,ALTER,ADD,COMMIT,UNSIGNED}
}
\lstset{style=sql}

\tcbset{
  sqlbox/.style={colback=codbg,colframe=sqlpurple!55,
    fonttitle=\bfseries\small\color{white},coltitle=white,
    attach boxed title to top left={yshift=-2mm,xshift=4mm},
    boxed title style={colback=sqlpurple,colframe=sqlpurple,arc=2pt},
    boxrule=0.5pt,arc=3pt,left=4pt,right=4pt,top=6pt,bottom=4pt,breakable},
  resbox/.style={colback=resbg,colframe=sqlgreen!55,
    fonttitle=\bfseries\small\color{white},coltitle=white,
    attach boxed title to top left={yshift=-2mm,xshift=4mm},
    boxed title style={colback=sqlgreen!80!black,colframe=sqlgreen,arc=2pt},
    boxrule=0.5pt,arc=3pt,left=4pt,right=4pt,top=6pt,bottom=4pt},
  notebox/.style={colback=notecolor,colframe=sqlorange!60,
    fonttitle=\bfseries\small\color{white},coltitle=white,
    attach boxed title to top left={yshift=-2mm,xshift=4mm},
    boxed title style={colback=sqlorange,colframe=sqlorange,arc=2pt},
    boxrule=0.5pt,arc=3pt,left=4pt,right=4pt,top=6pt,bottom=4pt,breakable},
  warnbox/.style={colback=warncolor,colframe=sqlred!55,
    fonttitle=\bfseries\small\color{white},coltitle=white,
    attach boxed title to top left={yshift=-2mm,xshift=4mm},
    boxed title style={colback=sqlred,colframe=sqlred,arc=2pt},
    boxrule=0.5pt,arc=3pt,left=4pt,right=4pt,top=6pt,bottom=4pt,breakable}
}

\titleformat{\section}{\large\bfseries\color{sqlpurple}}{}{0em}{}[\color{sqlpurple!40}\titlerule]
\titleformat{\subsection}{\small\bfseries\color{sqlblue}}{}{0em}{$\bullet$\ }
\titlespacing*{\section}{0pt}{8pt}{4pt}
\titlespacing*{\subsection}{0pt}{5pt}{2pt}
\setlist[itemize]{noitemsep,topsep=1pt,leftmargin=*}

\newcommand{\TBLH}[1]{\textbf{\textcolor{white}{#1}}}

\begin{document}

%% TITEL
\begin{tcolorbox}[colback=sqlpurple!10,colframe=sqlpurple,boxrule=1.2pt,arc=5pt,top=4pt,bottom=4pt]
\begin{center}
{\Large\bfseries\color{sqlpurple} SQL Samenvatting --- IEE Databank}\\[1pt]
{\small Databeheer ABL 2025--2026 \quad|\quad Cheat sheet mag gebruikt worden op examen}
\end{center}
\end{tcolorbox}
\vspace{4pt}

%%=============================================================
\section{De IEE Databank --- Structuur}
%%=============================================================

\begin{multicols}{3}
{\scriptsize
\begin{tabular}{|l|l|l|}
\hline
\rowcolor{tblhead}\TBLH{employee} & \TBLH{type} & \TBLH{} \\ \hline
\underline{employee\_nr} & INT & PK \\ \hline
name & VARCHAR & \\ \hline
firstname & VARCHAR & \\ \hline
func & VARCHAR & \\ \hline
\textit{supervisor} & INT & FK* \\ \hline
birthdate & DATE & \\ \hline
salary & DOUBLE & \\ \hline
\textit{department\_nr} & INT & FK \\ \hline
\multicolumn{3}{|l|}{\tiny *FK naar employee zelf} \\ \hline
\end{tabular}}

\vspace{3pt}

{\scriptsize
\begin{tabular}{|l|l|l|}
\hline
\rowcolor{tblhead}\TBLH{department} & \TBLH{type} & \TBLH{} \\ \hline
\underline{department\_nr} & INT & PK \\ \hline
name & VARCHAR & \\ \hline
location & VARCHAR & \\ \hline
\textit{head} & INT & FK \\ \hline
\multicolumn{3}{|l|}{\tiny head = FK naar employee} \\ \hline
\end{tabular}}

\columnbreak

{\scriptsize
\begin{tabular}{|l|l|l|}
\hline
\rowcolor{tblhead}\TBLH{course} & \TBLH{type} & \TBLH{} \\ \hline
\underline{code} & CHAR(4) & PK \\ \hline
description & VARCHAR & \\ \hline
type & VARCHAR(4) & \\ \hline
duration & INT & \\ \hline
\multicolumn{3}{|l|}{\tiny types: BLD, ALG, DSG} \\ \hline
\end{tabular}}

\vspace{3pt}

{\scriptsize
\begin{tabular}{|l|l|l|}
\hline
\rowcolor{tblhead}\TBLH{execution} & \TBLH{type} & \TBLH{} \\ \hline
\underline{\textit{course}} & CHAR(4) & PK+FK \\ \hline
\underline{begindate} & DATETIME & PK \\ \hline
\textit{teacher} & INT & FK \\ \hline
location & VARCHAR & \\ \hline
\multicolumn{3}{|l|}{\tiny PK = course+begindate samen!} \\ \hline
\end{tabular}}

\columnbreak

{\scriptsize
\begin{tabular}{|l|l|l|}
\hline
\rowcolor{tblhead}\TBLH{subscription} & \TBLH{type} & \TBLH{} \\ \hline
\underline{sub\_id} & INT & PK \\ \hline
\textit{student} & INT & FK \\ \hline
\textit{course} & CHAR(4) & FK \\ \hline
\textit{begindate} & DATETIME & FK \\ \hline
evaluation & INT & \\ \hline
\multicolumn{3}{|l|}{\tiny student = FK naar employee!} \\ \hline
\multicolumn{3}{|l|}{\tiny course+begindate = FK naar execution} \\ \hline
\end{tabular}}
\end{multicols}

%%=============================================================
\section{Data in de Tabellen}
%%=============================================================

\subsection{employee}

{\scriptsize
\begin{tabular}{|c|l|l|l|c|l|r|c|}
\hline
\rowcolor{tblhead}\TBLH{emp\_nr} & \TBLH{name} & \TBLH{firstname} & \TBLH{func} & \TBLH{supervisor} & \TBLH{birthdate} & \TBLH{salary} & \TBLH{dept} \\ \hline
7369 & Smit     & Nick       & trainer     & 7902 & 1964-02-15 & 1979.92 & 20 \\ \hline
\rowcolor{rowalt}
7499 & Alders   & Jan        & salesperson & 7698 & 1957-02-18 & 2259.84 & 30 \\ \hline
7521 & Degens   & Theo       & salesperson & 7698 & 1978-02-20 & 2049.88 & 30 \\ \hline
\rowcolor{rowalt}
7566 & Jansen   & Jean-Marie & manager     & 7839 & 1963-04-01 & 4284.71 & 20 \\ \hline
7654 & Martens  & Pierre     & salesperson & 7698 & 1952-09-27 & 2149.88 & 30 \\ \hline
\rowcolor{rowalt}
7698 & Bollen   & Roel       & manager     & 7839 & 1959-10-30 & 3109.72 & 30 \\ \hline
7782 & Clerckx  & Axelle     & manager     & 7839 & 1981-06-08 & 2869.76 & 10 \\ \hline
\rowcolor{rowalt}
7788 & Schotten & Sam        & trainer     & 7566 & 1955-11-24 & 2399.71 & 20 \\ \hline
7839 & Dekoning & Carlo      & CEO         & NULL & 1948-11-15 & 5299.51 & 10 \\ \hline
\rowcolor{rowalt}
7844 & Svenson  & Bjorn      & salesperson & 7698 & 1964-09-27 & 2199.85 & 30 \\ \hline
7876 & Li       & Muhan      & trainer     & 7788 & 1972-12-28 & 1959.89 & 20 \\ \hline
\rowcolor{rowalt}
7900 & Jansen   & Roel       & accountant  & 7698 & 1982-12-01 & 1879.92 & 30 \\ \hline
7902 & Swinnen  & Mark       & trainer     & 7566 & 1955-02-11 & 2299.71 & 20 \\ \hline
\rowcolor{rowalt}
7934 & Molenaar & Theo       & accountant  & 7782 & 1958-01-21 & 2179.87 & 10 \\ \hline
\end{tabular}}

\begin{multicols}{2}
\subsection{department}

{\scriptsize
\begin{tabular}{|c|l|l|c|}
\hline
\rowcolor{tblhead}\TBLH{nr} & \TBLH{name} & \TBLH{location} & \TBLH{head} \\ \hline
10 & headquarters    & Antwerpen & 7782 \\ \hline
\rowcolor{rowalt}
20 & education       & Leuven    & 7566 \\ \hline
30 & sales           & Kontich   & 7698 \\ \hline
\rowcolor{rowalt}
40 & human resources & Heverlee  & 7839 \\ \hline
50 & strategy        & Antwerpen & 7839 \\ \hline
\end{tabular}}

\columnbreak

\subsection{course}

{\scriptsize
\begin{tabular}{|l|l|c|c|}
\hline
\rowcolor{tblhead}\TBLH{code} & \TBLH{description (kort)} & \TBLH{type} & \TBLH{dur} \\ \hline
CGE & ORACLE CASE generation  & DSG & 4 \\ \hline
\rowcolor{rowalt}
CSO & ORACLE CASE design       & DSG & 5 \\ \hline
ERM & Data modeling ERM        & DSG & 3 \\ \hline
\rowcolor{rowalt}
FOR & SQL*Forms                & BLD & 4 \\ \hline
OAG & ORACLE app users         & ALG & 1 \\ \hline
\rowcolor{rowalt}
PMT & Process modeling         & DSG & 1 \\ \hline
RSO & Relational design        & DSG & 2 \\ \hline
\rowcolor{rowalt}
S02 & Intro SQL+SQL*Plus       & ALG & 4 \\ \hline
SMU & SQL*Menu                 & BLD & 1 \\ \hline
\rowcolor{rowalt}
SRW & SQL*ReportWriter         & BLD & 2 \\ \hline
\end{tabular}}
\end{multicols}

\begin{multicols}{2}
\subsection{execution}

{\scriptsize
\begin{tabular}{|l|l|c|l|}
\hline
\rowcolor{tblhead}\TBLH{course} & \TBLH{begindate} & \TBLH{teacher} & \TBLH{location} \\ \hline
CSO & 2018-02-16 & NULL & Leuven  \\ \hline
\rowcolor{rowalt}
ERM & 2018-01-12 & NULL & NULL    \\ \hline
FOR & 2017-10-14 & 7566 & Landen  \\ \hline
\rowcolor{rowalt}
FOR & 2018-02-09 & 7876 & Leuven  \\ \hline
OAG & 2017-08-10 & 7566 & Kontich \\ \hline
\rowcolor{rowalt}
OAG & 2017-09-28 & 7902 & Leuven  \\ \hline
RSO & 2018-02-16 & 7788 & Kontich \\ \hline
\rowcolor{rowalt}
S02 & 2017-03-10 & 7369 & Landen  \\ \hline
S02 & 2018-04-06 & 7369 & Leuven  \\ \hline
\rowcolor{rowalt}
S02 & 2018-04-21 & 7902 & Leuven  \\ \hline
SMU & 2017-09-11 & 7788 & Leuven  \\ \hline
\rowcolor{rowalt}
SRW & 2017-02-10 & 7369 & Leuven  \\ \hline
SRW & 2018-09-11 & NULL & Landen  \\ \hline
\end{tabular}}

\columnbreak

\subsection{subscription}

{\scriptsize
\begin{tabular}{|c|l|l|c|}
\hline
\rowcolor{tblhead}\TBLH{student} & \TBLH{course} & \TBLH{begindate} & \TBLH{eval} \\ \hline
7499 & FOR & 2017-10-14 & 4    \\ \hline
\rowcolor{rowalt}
7499 & S02 & 2018-04-21 & 8    \\ \hline
7499 & SMU & 2017-09-11 & 9    \\ \hline
\rowcolor{rowalt}
7499 & SRW & 2017-02-10 & 9    \\ \hline
7521 & OAG & 2017-08-10 & 7    \\ \hline
\rowcolor{rowalt}
7566 & FOR & 2018-02-09 & NULL \\ \hline
7566 & SMU & 2017-09-11 & 7    \\ \hline
\rowcolor{rowalt}
7698 & FOR & 2018-02-09 & NULL \\ \hline
7698 & S02 & 2018-04-06 & NULL \\ \hline
\rowcolor{rowalt}
7698 & S02 & 2018-04-21 & 8    \\ \hline
7782 & FOR & 2017-10-14 & 9    \\ \hline
\rowcolor{rowalt}
7788 & FOR & 2017-10-14 & 9    \\ \hline
7788 & FOR & 2018-02-09 & NULL \\ \hline
\rowcolor{rowalt}
7788 & S02 & 2017-03-10 & NULL \\ \hline
7839 & CSO & 2018-02-16 & NULL \\ \hline
\rowcolor{rowalt}
7839 & FOR & 2017-10-14 & 8    \\ \hline
7839 & RSO & 2018-02-16 & NULL \\ \hline
\rowcolor{rowalt}
7839 & S02 & 2017-03-10 & 6    \\ \hline
7844 & OAG & 2017-09-28 & 10   \\ \hline
\rowcolor{rowalt}
7876 & FOR & 2017-10-14 & 10   \\ \hline
7876 & S02 & 2018-04-21 & 10   \\ \hline
\rowcolor{rowalt}
7876 & SMU & 2017-09-11 & 6    \\ \hline
7900 & OAG & 2017-08-10 & 7    \\ \hline
\rowcolor{rowalt}
7900 & SRW & 2017-02-10 & 4    \\ \hline
7902 & OAG & 2017-08-10 & 9    \\ \hline
\rowcolor{rowalt}
7902 & S02 & 2017-03-10 & 8    \\ \hline
7902 & S02 & 2018-04-06 & NULL \\ \hline
\rowcolor{rowalt}
7934 & S02 & 2018-04-21 & 9    \\ \hline
\end{tabular}}
\end{multicols}

%%=============================================================
\section{DDL --- Tabellen Aanmaken}
%%=============================================================

\begin{tcolorbox}[sqlbox,title=Structuur + constraints]
\begin{lstlisting}
-- Samengestelde PRIMARY KEY (course + begindate SAMEN uniek):
CREATE TABLE execution (
    course    CHAR(4),
    begindate DATETIME,
    teacher   INTEGER,
    location  VARCHAR(50),
    PRIMARY KEY (course, begindate),    -- beide samen = PK
    FOREIGN KEY (course)  REFERENCES course(code),
    FOREIGN KEY (teacher) REFERENCES employee(employee_nr)
);

-- Zelfrelatie: supervisor FK verwijst naar employee zelf
CREATE TABLE employee (
    employee_nr   INTEGER PRIMARY KEY,
    supervisor    INTEGER,
    department_nr INTEGER,
    FOREIGN KEY (supervisor)    REFERENCES employee(employee_nr),
    FOREIGN KEY (department_nr) REFERENCES department(department_nr)
);
\end{lstlisting}
\end{tcolorbox}

\begin{tcolorbox}[notebox,title=Constraints samengevat]
{\footnotesize\begin{tabular}{ll}
\texttt{\textcolor{sqlpurple}{\bfseries NOT NULL}} & Kolom mag niet leeg zijn \\
\texttt{\textcolor{sqlpurple}{\bfseries UNIQUE}} & Waarde uniek in de tabel \\
\texttt{\textcolor{sqlpurple}{\bfseries DEFAULT x}} & Standaardwaarde als niets opgegeven \\
\texttt{\textcolor{sqlpurple}{\bfseries AUTO\_INCREMENT}} & Automatisch ophogen (voor PK) \\
\texttt{\textcolor{sqlpurple}{\bfseries PRIMARY KEY (a,b)}} & Combinatie a+b is uniek \\
\texttt{\textcolor{sqlpurple}{\bfseries FOREIGN KEY ... REFERENCES}} & Link naar PK van andere tabel \\
\end{tabular}}
\end{tcolorbox}

%%=============================================================
\section{Basis SELECT}
%%=============================================================

\textbf{Uitvoervolgorde:} FROM(1) $\to$ WHERE(2) $\to$ GROUP BY(3) $\to$ HAVING(4) $\to$ SELECT(5) $\to$ ORDER BY(6)

\begin{tcolorbox}[sqlbox,title=Basisvoorbeelden]
\begin{lstlisting}
SELECT * FROM employee;
SELECT employee_nr, salary*12 AS yearly_salary FROM employee;
SELECT DISTINCT func FROM employee;

-- Filteren
SELECT name, firstname FROM employee WHERE birthdate >= '19600101';
SELECT * FROM employee WHERE NOT (name = 'Jansen' AND firstname = 'Roel');
-- Alternatief: WHERE name != 'Jansen' OR firstname != 'Roel'

-- NULL checken
SELECT * FROM subscription WHERE evaluation IS NULL;
SELECT * FROM subscription WHERE evaluation IS NOT NULL;

-- NOT IN: let op NULL! (zie valkuilen)
SELECT * FROM subscription WHERE evaluation NOT IN (7,8,9);
SELECT * FROM subscription WHERE evaluation NOT IN (7,8,9) OR evaluation IS NULL;

-- Tekst samenvoegen + sorteren
SELECT CONCAT(LEFT(firstname,1),'.', name) FROM employee;
SELECT department_nr, name FROM employee ORDER BY department_nr, name ASC;
\end{lstlisting}
\end{tcolorbox}

\begin{tcolorbox}[notebox,title=Wildcards bij LIKE]
\texttt{\%} = 0 of meer tekens \quad|\quad \texttt{\_} = exact 1 teken\\
\texttt{LIKE 'J\%'} $\to$ begint met J \quad|\quad
\texttt{LIKE '\%en'} $\to$ eindigt op en \quad|\quad
\texttt{LIKE '\_\_n\%'} $\to$ 3e letter is n
\end{tcolorbox}

%%=============================================================
\section{GROUP BY \& HAVING}
%%=============================================================

\begin{tcolorbox}[sqlbox,title=Voorbeeld 1: count + avg per department]
\begin{lstlisting}
SELECT department_nr, COUNT(*) AS number_of_employees, AVG(salary) AS average
FROM employee
GROUP BY department_nr;
\end{lstlisting}
\end{tcolorbox}

\begin{tcolorbox}[resbox,title=Resultaat]
{\scriptsize\begin{tabular}{|c|c|c|}
\hline
\rowcolor{tblhead}\TBLH{department\_nr} & \TBLH{number\_of\_employees} & \TBLH{average} \\ \hline
10 & 3 & 3449.71 \\ \hline
\rowcolor{rowalt}
20 & 5 & 2127.44 \\ \hline
30 & 5 & 2099.85 \\ \hline
\end{tabular}}
\end{tcolorbox}

\begin{tcolorbox}[sqlbox,title=Voorbeeld 2: HAVING --- studenten met gem evaluatie $>$ 7]
\begin{lstlisting}
SELECT student, COUNT(*) AS number_of_subscriptions
FROM subscription
GROUP BY student
HAVING AVG(evaluation) > 7;
\end{lstlisting}
\end{tcolorbox}

\begin{tcolorbox}[sqlbox,title=Voorbeeld 3: subscriptions per execution gesorteerd]
\begin{lstlisting}
SELECT course, begindate, COUNT(*) AS Number
FROM subscription
GROUP BY course, begindate   -- GROUP BY op 2 kolommen: execution = course+begindate!
ORDER BY Number DESC;
\end{lstlisting}
\end{tcolorbox}

\begin{tcolorbox}[notebox,title=WHERE vs HAVING]
{\footnotesize\begin{tabular}{lll}
 & \textbf{WHERE} & \textbf{HAVING} \\ \midrule
Wanneer? & Voor GROUP BY & Na GROUP BY \\
Filtert? & Individuele rijen & Groepen \\
Aggregaatfuncties (AVG, COUNT...)? & NEE & JA \\
\end{tabular}}\\[3pt]
\textbf{Regel:} alles in SELECT zonder aggregaatfunctie MOET in GROUP BY staan, anders error!
\end{tcolorbox}

%%=============================================================
\section{Subqueries}
%%=============================================================

\begin{tcolorbox}[sqlbox,title=Subquery met 1 resultaat: gebruik =]
\begin{lstlisting}
-- Employees met salary hoger dan het gemiddelde
SELECT * FROM employee
WHERE salary > (SELECT AVG(salary) FROM employee);

-- Employees ouder dan Martens Pierre
SELECT * FROM employee
WHERE birthdate > (
    SELECT birthdate FROM employee
    WHERE firstname = 'Pierre' AND name = 'Martens'
);
\end{lstlisting}
\end{tcolorbox}

\begin{tcolorbox}[sqlbox,title=Subquery met meerdere resultaten: gebruik IN]
\begin{lstlisting}
-- Namen van employees die een BLD-cursus gevolgd hebben (al begonnen)
SELECT name, firstname FROM employee
WHERE employee_nr IN (
    SELECT student FROM subscription
    WHERE course IN (
        SELECT code FROM course WHERE type = 'BLD'
    )
    AND begindate < NOW()
);
\end{lstlisting}
\end{tcolorbox}

\begin{tcolorbox}[resbox,title=Resultaat]
{\scriptsize\begin{tabular}{|l|l|}
\hline
\rowcolor{tblhead}\TBLH{name} & \TBLH{firstname} \\ \hline
Alders   & Jan        \\ \hline
\rowcolor{rowalt}
Bollen   & Roel       \\ \hline
Clerckx  & Axelle     \\ \hline
\rowcolor{rowalt}
Dekoning & Carlo      \\ \hline
Jansen   & Jean-Marie \\ \hline
\rowcolor{rowalt}
Li       & Muhan      \\ \hline
Schotten & Sam        \\ \hline
\rowcolor{rowalt}
Swinnen  & Mark       \\ \hline
\end{tabular}}\\[2pt]
\textit{student in subscription = employee\_nr! Er is geen aparte student-tabel.}
\end{tcolorbox}

%%=============================================================
\section{Gecorreleerde Subquery}
%%=============================================================

Voor elke rij in de buitenste query wordt de subquery opnieuw uitgevoerd met een waarde uit die rij. Verplicht een \textbf{alias} op de buitenste tabel!

\begin{tcolorbox}[sqlbox,title=Employees met salary $>$ gemiddelde van hun eigen department]
\begin{lstlisting}
SELECT department_nr, name, firstname, salary
FROM employee AS e
WHERE salary > (
    SELECT AVG(salary)
    FROM employee
    WHERE department_nr = e.department_nr   -- e. = waarde van buitenste rij!
)
ORDER BY department_nr;
\end{lstlisting}
\end{tcolorbox}

\begin{tcolorbox}[resbox,title=Resultaat]
{\scriptsize\begin{tabular}{|c|l|l|r|}
\hline
\rowcolor{tblhead}\TBLH{dept} & \TBLH{name} & \TBLH{firstname} & \TBLH{salary} \\ \hline
10 & Dekoning & Carlo      & 5299.51 \\ \hline
\rowcolor{rowalt}
20 & Jansen   & Jean-Marie & 4284.71 \\ \hline
20 & Schotten & Sam        & 2399.71 \\ \hline
\rowcolor{rowalt}
30 & Bollen   & Roel       & 3109.72 \\ \hline
30 & Alders   & Jan        & 2259.84 \\ \hline
\end{tabular}}\\[2pt]
\textit{Elke employee wordt vergeleken met het gemiddelde van zijn eigen dept, niet van alle.}
\end{tcolorbox}

%%=============================================================
\section{Joins}
%%=============================================================

Een JOIN plakt tabellen aan elkaar op basis van PK=FK. Als een employee 3 subscriptions heeft, verschijnt hij 3 keer.

\subsection{INNER JOIN --- alleen rijen met match in beide tabellen}

\begin{tcolorbox}[sqlbox,title=Employees + department naam]
\begin{lstlisting}
SELECT employee.name, func, department.name AS department_name
FROM employee
JOIN department ON employee.department_nr = department.department_nr
WHERE department.name != 'education'
ORDER BY location, employee.name;
\end{lstlisting}
\end{tcolorbox}

\begin{tcolorbox}[sqlbox,title=Zelfrelatie --- employee met zijn supervisor (alias verplicht!)]
\begin{lstlisting}
SELECT subordinate.name AS name_employee,
       superior.name    AS name_supervisor
FROM employee AS subordinate
JOIN employee AS superior ON subordinate.supervisor = superior.employee_nr;
-- Dekoning heeft supervisor=NULL -> verdwijnt bij INNER JOIN!
\end{lstlisting}
\end{tcolorbox}

\begin{tcolorbox}[resbox,title=Resultaat zelfrelatie INNER JOIN (deel)]
{\scriptsize\begin{tabular}{|l|l|}
\hline
\rowcolor{tblhead}\TBLH{name\_employee} & \TBLH{name\_supervisor} \\ \hline
Jansen (JM) & Dekoning \\ \hline
\rowcolor{rowalt}
Bollen      & Dekoning \\ \hline
Clerckx     & Dekoning \\ \hline
\rowcolor{rowalt}
Swinnen     & Jansen (JM) \\ \hline
Smit        & Swinnen \\ \hline
\end{tabular}}\\[2pt]
\textit{Dekoning (CEO, supervisor=NULL) verdwijnt bij INNER JOIN. Gebruik LEFT JOIN om hem te bewaren.}
\end{tcolorbox}

\begin{tcolorbox}[sqlbox,title=3 tabellen joinen --- employee + department head]
\begin{lstlisting}
SELECT subordinate.name AS employee,
       boss.name        AS department_head,
       department.name
FROM employee AS subordinate
JOIN department ON subordinate.department_nr = department.department_nr
JOIN employee AS boss ON department.head = boss.employee_nr
ORDER BY department.name;
\end{lstlisting}
\end{tcolorbox}

\begin{tcolorbox}[sqlbox,title=Grote join over 5 tabellen (uit oplossingen)]
\begin{lstlisting}
SELECT CONCAT(emp.name,' ',emp.firstname) AS employee,
       dep.name AS department,
       code, description,
       CONCAT(teach.name,' ',teach.firstname) AS teacher,
       subscription.begindate, evaluation
FROM employee AS emp
JOIN department AS dep ON emp.department_nr = dep.department_nr
JOIN subscription      ON subscription.student = emp.employee_nr
JOIN execution         ON execution.course = subscription.course
                      AND execution.begindate = subscription.begindate
JOIN employee AS teach ON teach.employee_nr = teacher
JOIN course            ON course.code = subscription.course
ORDER BY emp.name, emp.firstname, begindate;
\end{lstlisting}
\end{tcolorbox}

\subsection{LEFT JOIN --- alle rijen van linkertabel, ook zonder match}

\begin{tcolorbox}[sqlbox,title=CEO bewaren die geen supervisor heeft]
\begin{lstlisting}
SELECT subordinate.name AS name_employee,
       superior.name    AS name_supervisor
FROM employee AS subordinate
LEFT JOIN employee AS superior ON subordinate.supervisor = superior.employee_nr;
-- Dekoning: supervisor=NULL -> blijft in resultaat, superior.name = NULL
\end{lstlisting}
\end{tcolorbox}

\begin{tcolorbox}[resbox,title=Verschil INNER vs LEFT JOIN voor Dekoning]
{\scriptsize\begin{tabular}{|l|c|c|}
\hline
\rowcolor{tblhead}\TBLH{name\_employee} & \TBLH{INNER JOIN} & \TBLH{LEFT JOIN} \\ \hline
Dekoning    & verdwijnt & NULL (blijft!) \\ \hline
\rowcolor{rowalt}
Jansen (JM) & Dekoning  & Dekoning \\ \hline
Bollen      & Dekoning  & Dekoning \\ \hline
\end{tabular}}
\end{tcolorbox}

\begin{tcolorbox}[notebox,title=INNER vs LEFT samengevat]
{\footnotesize\begin{tabular}{lcc}
\textbf{Situatie} & \textbf{INNER JOIN} & \textbf{LEFT JOIN} \\ \midrule
Match in beide tabellen & tonen & tonen \\
Alleen in linkertabel & verdwijnt & tonen met NULL \\
Alleen in rechtertabel & verdwijnt & verdwijnt \\
\end{tabular}}\\[2pt]
Gebruik LEFT JOIN als je rijen van de linkertabel \textit{altijd} wil tonen, ook zonder match.
\end{tcolorbox}

%%=============================================================
\section{Valkuilen --- Veelgemaakte Fouten}
%%=============================================================

\begin{tcolorbox}[warnbox,title=FOUT 1: = gebruiken als subquery meerdere rijen geeft]
\begin{lstlisting}
-- ERROR: Subquery returns more than 1 row
WHERE employee_nr = (SELECT student FROM subscription);
-- subscription heeft 28 rijen -> = weet niet welke te kiezen!

-- OPLOSSING: gebruik IN
WHERE employee_nr IN (SELECT student FROM subscription);

-- OOK FOUT: GROUP BY in subquery geeft meerdere rijen
WHERE salary > (SELECT AVG(salary) FROM employee GROUP BY department_nr);
-- Geeft 1 AVG per department -> meerdere rijen -> ERROR!
-- OPLOSSING: gecorreleerde subquery (zie sectie 6)
\end{lstlisting}
\end{tcolorbox}

\begin{tcolorbox}[warnbox,title=FOUT 2: Gecorreleerde subquery zonder alias]
\begin{lstlisting}
-- FOUT: vergelijkt department_nr met ZICHZELF -> altijd TRUE!
WHERE salary > (SELECT AVG(salary) FROM employee
                WHERE employee.department_nr = department_nr);
-- employee.department_nr EN department_nr zijn DEZELFDE kolom -> altijd gelijk!
-- Bovendien: geen WHERE-filter -> GROUP BY nodig -> meerdere rijen -> ERROR

-- OPLOSSING: alias e. op de buitenste tabel zodat je onderscheid maakt
FROM employee AS e
WHERE salary > (SELECT AVG(salary) FROM employee
                WHERE department_nr = e.department_nr);
-- e.department_nr = waarde van de huidige rij in de buitenste query
-- department_nr   = kolom van de binnenste tabel employee
\end{lstlisting}
\end{tcolorbox}

\begin{tcolorbox}[warnbox,title=FOUT 3: NOT IN faalt bij NULLs in de subquery]
\begin{lstlisting}
-- FOUT: execution.location heeft NULL-waarden (ERM en SRW 2018)
-- NOT IN (..., NULL) geeft altijd 0 rijen terug!
SELECT location FROM department
WHERE location NOT IN (SELECT location FROM execution);

-- OPLOSSING: filter NULLs weg uit de subquery
SELECT DISTINCT location FROM department
WHERE location NOT IN (
    SELECT location FROM execution WHERE location IS NOT NULL
);
\end{lstlisting}
\textbf{Waarom?} \texttt{NOT IN (Leuven, Kontich, NULL)} is UNKNOWN voor elke waarde.
SQL kan niet zeker zeggen of iets gelijk is aan NULL, dus gooit het de rij weg.\\[2pt]
\textit{Resultaat na fix: Heverlee en Antwerpen (dept 40 en 50 hebben geen cursussen)}
\end{tcolorbox}

\begin{tcolorbox}[warnbox,title=FOUT 4: Samengestelde FK checken met 1 kolom]
\begin{lstlisting}
-- FOUT: checkt alleen datum, NIET welke cursus op die datum!
AND begindate IN (SELECT begindate FROM execution)
-- Als FOR en OAG op dezelfde datum staan, matcht FOR ook rijen van OAG

-- OPLOSSING: check course EN begindate samen als combinatie
AND (course, begindate) IN (
    SELECT course, begindate FROM execution WHERE begindate <= NOW()
);
\end{lstlisting}
\end{tcolorbox}

%%=============================================================
\section{Snelle Referentie}
%%=============================================================

\begin{multicols}{2}
{\footnotesize\begin{tabular}{ll}
\toprule
\multicolumn{2}{l}{\textbf{Operatoren}} \\ \midrule
\texttt{=, <>, <, >, <=, >=} & Vergelijken \\
\texttt{AND, OR, NOT} & Logisch \\
\texttt{IN (...)} & In lijst \\
\texttt{NOT IN (...)} & Niet in lijst (let op NULL!) \\
\texttt{BETWEEN x AND y} & Bereik (inclusief) \\
\texttt{LIKE 'x\%'} & Begint met x \\
\texttt{LIKE '\%x'} & Eindigt op x \\
\texttt{LIKE '\_x\%'} & 2e letter is x \\
\texttt{IS NULL} & Leeg \\
\texttt{IS NOT NULL} & Niet leeg \\
\bottomrule
\end{tabular}}

\columnbreak

{\footnotesize\begin{tabular}{ll}
\toprule
\multicolumn{2}{l}{\textbf{Functies \& keywords}} \\ \midrule
\texttt{COUNT(*)} & Aantal rijen \\
\texttt{SUM(k)} & Som \\
\texttt{AVG(k)} & Gemiddelde \\
\texttt{MIN(k) / MAX(k)} & Min / Max \\
\texttt{DISTINCT} & Geen duplicaten \\
\texttt{AS naam} & Alias geven \\
\texttt{ASC / DESC} & Volgorde (ASC=standaard) \\
\texttt{NOW()} & Huidige datum+tijd \\
\texttt{CONCAT(a,b)} & Teksten samenvoegen \\
\texttt{LEFT(str,n)} & Eerste n tekens \\
\bottomrule
\end{tabular}}
\end{multicols}

\begin{tcolorbox}[notebox,title=Algemeen formaat --- schrijfvolgorde]
\begin{lstlisting}
SELECT   kolommen of functies       -- stap 5: wat tonen
FROM     tabel(len)                 -- stap 1: welke tabel
WHERE    conditie op rijen          -- stap 2: filter voor groepering
GROUP BY kolommen                   -- stap 3: maak groepjes
HAVING   conditie op groepen        -- stap 4: filter na groepering
ORDER BY kolommen ASC|DESC;         -- stap 6: sorteer
\end{lstlisting}
\end{tcolorbox}

\end{document}